\documentclass[12pt,a4paper]{exam}
\usepackage[utf8]{inputenc}
\usepackage{amsmath,amssymb,amsfonts}
\usepackage{geometry}
\usepackage{enumitem}
\usepackage{graphicx}
\usepackage{hyperref}
\usepackage{xcolor}
\geometry{a4paper, top=2cm, bottom=2cm, left=2.5cm, right=2.5cm}
\hypersetup{colorlinks=true, linkcolor=blue, urlcolor=blue}
\renewcommand{\thequestion}{\textbf{\arabic{question}}}
\renewcommand{\questionlabel}{\thequestion.}
\pointname{ marks}
\pointsdroppedatright
\bracketedpoints
\setlength{\parindent}{0pt}
\setlength{\emergencystretch}{3em}
\allowdisplaybreaks
\newsavebox{\schmathbox}
\newcommand{\fitmath}[1]{%
  \sbox{\schmathbox}{$\displaystyle #1$}%
  \ifdim\wd\schmathbox>\linewidth
    \resizebox{\linewidth}{!}{\usebox{\schmathbox}}%
  \else
    \usebox{\schmathbox}%
  \fi}

\begin{document}

\thispagestyle{empty}
\begin{center}
  \textbf{\LARGE Diag} \\[0.3cm]
  \textbf{Time Allowed: 3 Hours} \hfill \textbf{Maximum Marks: 80} \\[0.3cm]
  \rule{\textwidth}{0.5pt}
\end{center}

\vspace{0.3cm}

\noindent\textbf{General Instructions:}
\begin{enumerate}
  \item {\normalfont\normalcolor This question paper contains 40 questions. All questions are compulsory.}
  \item {\normalfont\normalcolor Questions are grouped into sections A through E.}
  \item {\normalfont\normalcolor Use of calculators is NOT permitted.}
\end{enumerate}
\bigskip
\noindent\textbf{\large Section A: Multiple Choice \& Assertion-Reason (1 mark each)}

\begin{questions}
\question[1] Assertion (A): The length of a tangent drawn from an external point to a circle is always equal to the radius of the circle. Reason (R): A tangent to a circle is perpendicular to the radius at the point of contact.
\begin{enumerate}[label=(\Alph*)]
  \item (A) Both Assertion (A) and Reason (R) are true and Reason (R) is the correct explanation of Assertion (A)
  \item (B) Both Assertion (A) and Reason (R) are true but Reason (R) is NOT the correct explanation of Assertion (A)
  \item (C) Assertion (A) is true but Reason (R) is false
  \item (D) Assertion (A) is false but Reason (R) is true
\end{enumerate}
\vspace{0.15cm}

\question[1] Assertion (A): The distance between two parallel tangents of a circle of radius $3$ $cm$ is $6$ $cm$. Reason (R): The distance between two parallel tangents is equal to the diameter of the circle.
\begin{enumerate}[label=(\Alph*)]
  \item (A) Both Assertion (A) and Reason (R) are true and Reason (R) is the correct explanation of Assertion (A)
  \item (B) Both Assertion (A) and Reason (R) are true but Reason (R) is NOT the correct explanation of Assertion (A)
  \item (C) Assertion (A) is true but Reason (R) is false
  \item (D) Assertion (A) is false but Reason (R) is true
\end{enumerate}
\vspace{0.15cm}

\question[1] Assertion (A): If the angle between two tangents drawn from a point $P$ to a circle of radius $a$ and centre $O$ is $60^\circ$, then $OP = 2a$. Reason (R): The line joining the centre to the external point bisects the angle between the tangents.

\begin{center}
\includegraphics[width=0.48\textwidth]{/home/user/Desktop/backup/schEngine/pdfs/diag/diagrams/diagram_009.pdf}
\end{center}

\begin{enumerate}[label=(\Alph*)]
  \item (A) Both Assertion (A) and Reason (R) are true and Reason (R) is the correct explanation of Assertion (A)
  \item (B) Both Assertion (A) and Reason (R) are true but Reason (R) is NOT the correct explanation of Assertion (A)
  \item (C) Assertion (A) is true but Reason (R) is false
  \item (D) Assertion (A) is false but Reason (R) is true
\end{enumerate}
\vspace{0.15cm}

\question[1] Assertion (A): The number of tangents that can be drawn from a point inside a circle is zero. Reason (R): A tangent is a line that intersects the circle at only one point.
\begin{enumerate}[label=(\Alph*)]
  \item (A) Both Assertion (A) and Reason (R) are true and Reason (R) is the correct explanation of Assertion (A)
  \item (B) Both Assertion (A) and Reason (R) are true but Reason (R) is NOT the correct explanation of Assertion (A)
  \item (C) Assertion (A) is true but Reason (R) is false
  \item (D) Assertion (A) is false but Reason (R) is true
\end{enumerate}
\vspace{0.15cm}

\question[1] Assertion (A): The tangents drawn at the ends of a diameter of a circle are parallel. Reason (R): The angles between the diameter and the tangents at the ends are $90^\circ$, and since the sum of interior angles on the same side of the transversal is $180^\circ$, they are parallel.
\begin{enumerate}[label=(\Alph*)]
  \item (A) Both Assertion (A) and Reason (R) are true and Reason (R) is the correct explanation of Assertion (A)
  \item (B) Both Assertion (A) and Reason (R) are true but Reason (R) is NOT the correct explanation of Assertion (A)
  \item (C) Assertion (A) is true but Reason (R) is false
  \item (D) Assertion (A) is false but Reason (R) is true
\end{enumerate}
\vspace{0.15cm}

\question[1] Assertion (A): If in $\triangle ABC$ and $\triangle DEF$, $\frac{AB}{DE} = \frac{BC}{EF} = \frac{CA}{FD} = \frac{2}{3}$, then $\triangle ABC \sim \triangle DEF$. 
 Reason (R): Two triangles are similar if their corresponding sides are proportional (SSS similarity criterion).
\begin{enumerate}[label=(\Alph*)]
  \item (A) Both Assertion (A) and Reason (R) are true and Reason (R) is the correct explanation of Assertion (A)
  \item (B) Both Assertion (A) and Reason (R) are true but Reason (R) is NOT the correct explanation of Assertion (A)
  \item (C) Assertion (A) is true but Reason (R) is false
  \item (D) Assertion (A) is false but Reason (R) is true
\end{enumerate}
\vspace{0.15cm}

\question[1] Assertion (A): The length of a tangent drawn from an external point to a circle is always equal to the radius of the circle. Reason (R): The tangent at any point of a circle is perpendicular to the radius through the point of contact.
\begin{enumerate}[label=(\Alph*)]
  \item (A) Both Assertion (A) and Reason (R) are true and Reason (R) is the correct explanation of Assertion (A)
  \item (B) Both Assertion (A) and Reason (R) are true but Reason (R) is NOT the correct explanation of Assertion (A)
  \item (C) Assertion (A) is true but Reason (R) is false
  \item (D) Assertion (A) is false but Reason (R) is true
\end{enumerate}
\vspace{0.15cm}

\question[1] Assertion (A): If a quadrilateral $\text{ABCD}$ circumscribes a circle, then $AB + CD = AD + BC$. Reason (R): The lengths of tangents drawn from an external point to a circle are equal.
\begin{enumerate}[label=(\Alph*)]
  \item (A) Both Assertion (A) and Reason (R) are true and Reason (R) is the correct explanation of Assertion (A)
  \item (B) Both Assertion (A) and Reason (R) are true but Reason (R) is NOT the correct explanation of Assertion (A)
  \item (C) Assertion (A) is true but Reason (R) is false
  \item (D) Assertion (A) is false but Reason (R) is true
\end{enumerate}
\vspace{0.15cm}

\end{questions}
\bigskip
\noindent\textbf{\large Section B: Very Short Answer (2 marks each)}

\begin{questions}
\question[2] A quadrilateral $\text{ABCD}$ is drawn to circumscribe a circle. If $AB = 6$ cm, $BC = 7$ cm, and $CD = 4$ cm, find the length of $AD$.
\vspace{0.15cm}

\question[2] In the given figure, if $\angle OAB = 40^\circ$, find the measure of $\angle APB$, where $O$ is the center of the circle and $PA, PB$ are tangents from an external point $P$.

\begin{center}
\includegraphics[width=0.48\textwidth]{/home/user/Desktop/backup/schEngine/pdfs/diag/diagrams/diagram_012.pdf}
\end{center}

\vspace{0.15cm}

\question[2] A tangent $PQ$ at a point $P$ of a circle of radius $5$ $cm$ meets a line through the center $O$ at a point $Q$ so that $OQ = 12$ $cm$. Find the length of $PQ$.
\vspace{0.15cm}

\question[2] Two concentric circles are of radii $5$ $cm$ and $3$ $cm$. Find the length of the chord of the larger circle which touches the smaller circle.
\vspace{0.15cm}

\question[2] If tangents $PA$ and $PB$ from a point $P$ to a circle with center $O$ are inclined to each other at angle of $80^\circ$, then find $\angle POA$.
\vspace{0.15cm}

\question[2] A quadrilateral $\text{ABCD}$ is drawn to circumscribe a circle. If $AB = 6$ $cm$, $BC = 7$ $cm$ and $CD = 4$ $cm$, find the length of $AD$.
\vspace{0.15cm}

\question[2] In a circle of radius $5$ $cm$, $AB$ and $AC$ are two chords such that $AB = AC = 6$ $cm$. If the tangent at $A$ intersects the line $BC$ produced at $P$, find the length of $AP$.

\begin{center}
\includegraphics[width=0.48\textwidth]{/home/user/Desktop/backup/schEngine/pdfs/diag/diagrams/diagram_013.pdf}
\end{center}

\vspace{0.15cm}

\question[2] If $\angle ATO = 40^\circ$ where $T$ is a point outside the circle with center $O$ and $TA$, $TB$ are tangents, find $\angle AOB$.

\begin{center}
\includegraphics[width=0.48\textwidth]{/home/user/Desktop/backup/schEngine/pdfs/diag/diagrams/diagram_014.pdf}
\end{center}

\vspace{0.15cm}

\question[2] Two concentric circles are of radii $5$ $cm$ and $3$ $cm$. Find the length of the chord of the larger circle which touches the smaller circle.
\vspace{0.15cm}

\question[2] If a quadrilateral $\text{ABCD}$ circumscribes a circle, prove that $AB + CD = AD + BC$.
\vspace{0.15cm}

\question[2] A point $P$ is $13$ $cm$ from the center of a circle. The length of the tangent drawn from $P$ to the circle is $12$ $cm$. Find the radius of the circle.

\begin{center}
\includegraphics[width=0.48\textwidth]{/home/user/Desktop/backup/schEngine/pdfs/diag/diagrams/diagram_016.pdf}
\end{center}

\vspace{0.15cm}

\end{questions}
\bigskip
\noindent\textbf{\large Section C: Short Answer (3 marks each)}

\begin{questions}
\question[3] A tangent $PQ$ at a point $P$ of a circle of radius $5$ $cm$ meets a line through the centre $O$ at a point $Q$ so that $OQ = 13$ $cm$. Find the length of $PQ$.

\begin{center}
\includegraphics[width=0.48\textwidth]{/home/user/Desktop/backup/schEngine/pdfs/diag/diagrams/diagram_017.pdf}
\end{center}

\vspace{0.15cm}

\question[3] Two concentric circles are of radii $5$ $cm$ and $3$ $cm$. Find the length of the chord of the larger circle which touches the smaller circle.
\vspace{0.15cm}

\question[3] If tangents $PA$ and $PB$ from a point $P$ to a circle with centre $O$ are inclined to each other at an angle of $80^\circ$, find $\angle POA$.
\vspace{0.15cm}

\question[3] A quadrilateral $\text{ABCD}$ is drawn to circumscribe a circle. Prove that $AB + CD = AD + BC$.
\vspace{0.15cm}

\question[3] Prove that the perpendicular at the point of contact to the tangent to a circle passes through the centre of the circle.
\vspace{0.15cm}

\question[3] Two concentric circles are of radii $5$ cm and $3$ cm. Find the length of the chord of the larger circle which touches the smaller circle.
\vspace{0.15cm}

\question[3] If a quadrilateral $\text{ABCD}$ is drawn to circumscribe a circle, prove that $AB + CD = AD + BC$.
\vspace{0.15cm}

\question[3] A tangent $PQ$ at a point $P$ of a circle of radius $5$ cm meets a line through the center $O$ at a point $Q$ so that $OQ = 12$ cm. Find the length of $PQ$.
\vspace{0.15cm}

\question[3] In $\triangle ABC$, a circle is inscribed touching sides $AB$, $BC$, and $CA$ at $D$, $E$, and $F$ respectively. If $AB = 12$ cm, $BC = 8$ cm, and $AC = 10$ cm, find the lengths of $AD$, $BE$, and $CF$.
\vspace{0.15cm}

\question[3] Two tangents $TP$ and $TQ$ are drawn to a circle with center $O$ from an external point $T$. Prove that $\angle PTQ = 2\angle OPQ$.
\vspace{0.15cm}

\question[3] In $\triangle ABC$, $DE \parallel BC$ such that $D$ lies on $AB$ and $E$ lies on $AC$. If $AD = 2.4 \text{ cm}$, $AE = 3.2 \text{ cm}$, and $EC = 4.8 \text{ cm}$, find the length of $AB$.

\begin{center}
\includegraphics[width=0.48\textwidth]{/home/user/Desktop/backup/schEngine/pdfs/diag/diagrams/diagram_019.pdf}
\end{center}

\vspace{0.15cm}

\end{questions}
\bigskip
\noindent\textbf{\large Section D: Long Answer (4-5 marks each)}

\begin{questions}
\question[5] Prove that the lengths of tangents drawn from an external point to a circle are equal. Using this result, if a quadrilateral $\text{ABCD}$ is drawn to circumscribe a circle, prove that $AB + CD = AD + BC$.

\begin{center}
\includegraphics[width=0.48\textwidth]{/home/user/Desktop/backup/schEngine/pdfs/diag/diagrams/diagram_020.pdf}
\end{center}

\vspace{0.15cm}

\question[5] Two concentric circles are of radii $5 \text{ cm}$ and $3 \text{ cm}$. Find the length of the chord of the larger circle which touches the smaller circle.

\begin{center}
\includegraphics[width=0.48\textwidth]{/home/user/Desktop/backup/schEngine/pdfs/diag/diagrams/diagram_021.pdf}
\end{center}

\vspace{0.15cm}

\question[5] A circle is inscribed in a $\triangle ABC$ having sides $8 \text{ cm}, 10 \text{ cm}$ and $12 \text{ cm}$. Find the lengths of the segments $AD, BE, CF$ where $D, E, F$ are points of contact on sides $AB, BC, CA$ respectively.
\vspace{0.15cm}

\question[5] Prove that the tangent at any point of a circle is perpendicular to the radius through the point of contact.

\begin{center}
\includegraphics[width=0.48\textwidth]{/home/user/Desktop/backup/schEngine/pdfs/diag/diagrams/diagram_022.pdf}
\end{center}

\vspace{0.15cm}

\question[5] Prove that the lengths of tangents drawn from an external point to a circle are equal. Using this, if a quadrilateral $\text{ABCD}$ is drawn to circumscribe a circle, prove that $AB + CD = AD + BC$.

\begin{center}
\includegraphics[width=0.48\textwidth]{/home/user/Desktop/backup/schEngine/pdfs/diag/diagrams/diagram_023.pdf}
\end{center}

\vspace{0.15cm}

\question[5] Two concentric circles are of radii $5 \ cm$ and $3 \ cm$. Find the length of the chord of the larger circle which touches the smaller circle.

\begin{center}
\includegraphics[width=0.48\textwidth]{/home/user/Desktop/backup/schEngine/pdfs/diag/diagrams/diagram_024.pdf}
\end{center}

\vspace{0.15cm}

\question[5] In $\triangle ABC$, a circle is inscribed touching $BC, CA,$ and $AB$ at $D, E,$ and $F$ respectively. If $AB = 12 \ cm, BC = 8 \ cm,$ and $AC = 10 \ cm$, find the lengths $AD, BE,$ and $CF$.
\vspace{0.15cm}

\question[5] Prove that the angle between the two tangents drawn from an external point to a circle is supplementary to the angle subtended by the line segment joining the points of contact at the centre.

\begin{center}
\includegraphics[width=0.48\textwidth]{/home/user/Desktop/backup/schEngine/pdfs/diag/diagrams/diagram_025.pdf}
\end{center}

\vspace{0.15cm}

\question[5] Prove that the lengths of tangents drawn from an external point to a circle are equal.

\begin{center}
\includegraphics[width=0.48\textwidth]{/home/user/Desktop/backup/schEngine/pdfs/diag/diagrams/diagram_026.pdf}
\end{center}

\vspace{0.15cm}

\question[5] From a point $Q$, the length of the tangent to a circle is $24 \ cm$ and the distance of $Q$ from the centre is $25 \ cm$. Find the radius of the circle.

\begin{center}
\includegraphics[width=0.48\textwidth]{/home/user/Desktop/backup/schEngine/pdfs/diag/diagrams/diagram_027.pdf}
\end{center}

\vspace{0.15cm}

\end{questions}
\bigskip

\end{document}
